% ============================================================================
% APÊNDICE D — Referências e fontes de pesquisa
% ============================================================================
\chapter{Referências, Fontes de Pesquisa e Metodologia}
\label{apendice:referencias}

Este apêndice documenta \emph{como} o livro foi fundamentado: as fontes de
mercado, a prioridade dada à documentação oficial e as datas de consulta
(detalhadas no apêndice~\ref{apendice:manutencao}).

\section{Metodologia de pesquisa}

\begin{enumerate}[leftmargin=*, itemsep=0.4em]
  \item \textbf{Dados de mercado} foram extraídos de levantamentos públicos e
  independentes: Stack Overflow Developer Survey 2025, GitHub Octoverse 2025 e
  State of JavaScript 2024;
  \item \textbf{Conteúdo técnico} prioriza a documentação oficial de cada
  tecnologia (Next.js, React, TypeScript, PostgreSQL, Node.js, Tailwind,
  Prisma/Drizzle, Vitest, Playwright, Docker, AWS, OWASP);
  \item \textbf{Revisão por pares simulada} (apêndice~\ref{apendice:revisao})
  valida tecnicamente o conteúdo, cruza as fontes e verifica a pedagogia;
  \item Todo código do livro é executável no repositório \texttt{codigo/} e
  passa pelos testes da suíte de revisão.
\end{enumerate}

\section{Fontes primárias de mercado}

\begin{itemize}[leftmargin=*, itemsep=0.4em]
  \item Stack Overflow Developer Survey 2025 --- Technology: \url{https://survey.stackoverflow.co/2025/technology}
  \item Stack Overflow Developer Survey 2025 (geral): \url{https://survey.stackoverflow.co/2025}
  \item GitHub Octoverse 2025: \url{https://octoverse.github.com/}
  \item State of JavaScript 2024: \url{https://2024.stateofjs.com/}
  \item ThoughtWorks Technology Radar: \url{https://www.thoughtworks.com/radar}
\end{itemize}

\section{Documentação oficial por tecnologia}

\begin{itemize}[leftmargin=*, itemsep=0.4em]
  \item HTML/CSS/JS: MDN Web Docs --- \url{https://developer.mozilla.org}
  \item Next.js: \url{https://nextjs.org/docs}
  \item React: \url{https://react.dev}
  \item TypeScript: \url{https://www.typescriptlang.org/docs/}
  \item Node.js: \url{https://nodejs.org/docs/latest/api/}
  \item PostgreSQL 18: \url{https://www.postgresql.org/docs/current/}
  \item Tailwind CSS: \url{https://tailwindcss.com/docs}
  \item Prisma: \url{https://www.prisma.io/docs};\ \ Drizzle: \url{https://orm.drizzle.team}
  \item Vitest: \url{https://vitest.dev};\ \ Playwright: \url{https://playwright.dev}
  \item Docker: \url{https://docs.docker.com}
  \item GitHub Actions: \url{https://docs.github.com/actions}
  \item AWS: \url{https://docs.aws.amazon.com}
  \item OWASP Top 10:2025: \url{https://owasp.org/Top10/2025/en/}
  \item Vercel AI SDK: \url{https://ai-sdk.dev}
\end{itemize}

\section{Padrões de citação usados no livro}

\begin{itemize}[leftmargin=*, itemsep=0.4em]
  \item \textbf{Citação acadêmica (BibLaTeX)}: para dados de mercado e
  padrões (ex.: \cite{so2025}, \cite{octoverse2025}, \cite{owasp2025});
  \item \textbf{Referência direta à documentação oficial}: ao final de cada
  capítulo, na seção ``Referências oficiais para aprofundar'';
  \item \textbf{URLs no texto}: \texttt{\url{...}} quando a referência é
  operacional (ferramenta, validador, jogo educativo).
\end{itemize}

\begin{caixadica}
A bibliografia completa, em formato BibLaTeX, vive em
\texttt{livro/referencias.bib}. As datas de consulta e as versões registradas
de cada tecnologia estão no apêndice~\ref{apendice:manutencao} e no arquivo
\texttt{MANUTENCAO.md} — mantenha-os atualizados em cada revisão do livro.
\end{caixadica}
